\documentclass[11pt,a4paper]{article}

\usepackage[margin=2.5cm]{geometry}
\usepackage{enumitem}
\usepackage{hyperref}
\usepackage{parskip}
\usepackage{amsmath}

\setlist[itemize]{leftmargin=*,topsep=2pt,itemsep=1pt}

\newcommand{\sect}[1]{\medskip\noindent\textbf{#1}\par}
\newcommand{\q}[1]{``\textit{#1}''}
\newcommand{\pg}[1]{(p.~#1)}

\title{Opposition Review}
\author{Opponent: Paul Mayer}
\date{Seminar date: 2026-05-18}

\begin{document}

\maketitle

\section*{Opposition Information}
\begin{itemize}
  \item \textbf{Author:} Rishi Vijayvargiya
  \item \textbf{Title:} Stepwise Precision Refinement For Computing Symmetric Matrix Squares on Tensor Cores --- With Adaptive Splitting for Reduced Computation Redundancy
  \item \textbf{Opponent:} Paul Mayer
  \item \textbf{Date of opposition:} 2026-05-18
\end{itemize}

\section{Summary}

The thesis proposes \emph{stepwise precision refinement}, an extension of the FP32$\to$FP16 precision-refinement technique of Markidis et al.~to three working precisions (FP64$\to$FP32$\to$FP16), in order to approximate the square of an FP64 symmetric matrix on NVIDIA Tensor Cores. An FP64 input element is decomposed into four FP16 components ($A_1,\ldots,A_4$); the matrix square is then computed via 10 (instead of 16) Tensor-Core matmuls by exploiting symmetry. The author additionally proposes \emph{adaptive splitting}, a runtime filter that drops $A_i$ matrices whose non-zero fraction falls below a threshold (NZT). A worst-case error bound is derived under a simplified model, and experiments on an A100 measure $t_{\text{total}}$, $e_{\max}$, $e_{\text{avg}}$ versus \texttt{cublasDgemm} for matrices $n\in[256,16384]$ and four input ranges.

The main empirical finding is that using 3 FP16 matrices is a reasonable middle ground between accuracy and runtime, and that adaptive splitting effectively skips redundant computations when input magnitudes are small enough that residual matrices vanish in FP16.

\section{Title, Introduction, Background, Theory and Previous Work}

Title and abstract accurately describe the content.
Both stepwise refinement and adaptive splitting contributions are mentioned correctly.

The Introduction starts with an interesting, memorable hook; however, it may be slightly misleading.
The crash portrait in the Ariane 5 anecdote \pg{1} was an overflow bug, not really something to do with lower-precision arithmetic.

Chapter~2 is excessively long for a CS master's audience in my opinion.
Some Theory is introduced but never used downstream (or maybe I missed it?), e.g., subnormal number discussion \pg{9}.
Binary point representation is undergraduate textbook material, CUDA thread/block hierarchy, kernel launch lifecycle etc. could be condensed by roughly half.
The related work section motivates symmetric matrix squares (physics, finance, chemistry, graph theory) very well.
It might have been nice to use one of these applications to then validate the method, but that's a scope decision.
Coverage of other precision refinement is thorough and seems competent.
However, the Ozaki Scheme~II is dismissed as \q{quite intricate} \pg{25} rather than discussed substantively, which was disappointing.
Also, empirical comparison with the Ozaki scheme is missing, despite Section~3.4 explicitly identifying it as the alternative and motivating the work as filling \q{the apparent gap in the precision refinement techniques} \pg{27}.
The thesis only positions itself \emph{against} prior work conceptually, never numerically.

Research goals are stated as four bullet points \pg{3}, but no explicit research question.
The implicit question (\q{Can stepwise refinement to 3 precisions improve accuracy of symmetric-matrix squares acceptably vs.~FP64?}) can be inferred but is never clearly stated.

A highlight to me is Section 4.5 (Theory of worst-case error bound).
It is well-explained and validated experimentally.

\section{Method}
Justification is mostly \q{extends Markidis et al. from 2 to 3 precisions}.
The decision to route FP64$\to$FP32$\to$FP16 (rather than directly FP64$\to$FP16 with multi-term residual) is taken as given;
Prerequisites are stated in Section 4.5.
E.g. he explicitly lists assumptions (no underflow/overflow, $|\delta_{16}|\le u_{16}$, etc.).

The method is adequatly explained.
Figures 4.1--4.4 make the decomposition very clear, and Eqs.~(4.5)--(4.6) walk through the symmetric-matrix shortcut explicitly.
Implementation chapter (Ch.~5) is well-paced; CUDA listings are concrete and traceable.

\section{Results, Discussion, Conclusions}
All results are well described.
The author walks through each metric per range systematically.
The findings seem honest.
He openly reports that the approach is \emph{slower} than FP64 \texttt{cublasDgemm} for small matrices \pg{74}

Evaluation was done on single hardware platform.
No portability claims were tested.
As stated above, tests were performed on synthetic inputs only:
uniformly distributed random symmetric matrices across four arbitrary ranges.
None of the motivating applications (covariance matrices, density matrices, etc.) appear in evaluation.
Only 5 runs averaged \pg{69}; no variance/CI reported in most figures.
No comparison against Ozaki scheme implementations.

Some Limitations are clearly stated
The author honestly discusses memory consumption and coarseness of adaptive splitting.

Section 8.2 maps \q{1 FP16 matrix} and \q{2 FP16 matrices} onto existing approaches
(\q{can be thought of as the application of the existing precision refinement technique} \pg{92}), which is a clean way to compare.
However, no comparison vs.~Ozaki, vs.~direct FP64$\to$FP16 cast, or vs.~FP32 \texttt{cublasGemmEx} was performed, which would help place the work into a broader context.

\textbf{Are the concolusions credible and supported by the results?}
Yes, the recommendation \q{3 FP16 matrices is a good middle ground} is well-supported for the ranges tested.
Adaptive splitting is shown to help only when residuals are vanishingly small (Ranges 1--3); for Range 4 nothing is filtered (Table 7.1).
Generalisation beyond the four tested ranges is not established, it was out of scope to begin with.

\section{Figures, Tables, Bibliography}
Some Figures are particularly good, e.g. Figures 4.1, 4.7 or 5.1 were good additions to simplify the understanding of the work.
However, I dislike the use of figures for equations, e.g. Fig. 4.2, ..., 4.4.
Some Figures could have been code listings (e.g. Fig. 4.5, 4.6) but this may be too nitpicky.
Figure 2.11 contains a caption error: the caption says \q{Representing (78.5625) in FP16} but the figure actually shows (1024.5).

Listing 5.7 and 5.8 captions \pg{62}: both captioned \q{Combining the Products of 3 FP16 Matrices to Form Final FP32 Result},
but neither listing does that (copy paste error).
I \emph{think} 5.7 is the overall algorithm, 5.8 is the FP64 reference \texttt{cublasDgemm} call. Fix that.

Figures that display results are done very well.
Axis and scales are well picked, transparent and readable.
They strongly support the result section.

I liked Table 7.5 \pg{88}: The colour-coded digit differences is a clever way to surface differences that the bar charts cannot show.
Figures 7.5/7.6 (line + bar duplication for execution time) feels redundant --- the appendix contains essentially the same plots for Ranges 2--3, suggesting a single combined figure would have sufficed.

What is missing?
A plot of the \emph{accuracy/runtime trade-off frontier} (e.g.,~$e_{\text{avg}}$ vs.~$t_{\text{total}}$ across configurations).
Could have made the main recommendation visually clear.

All adapted figures are credited.

The bibliography seems appropriate.
Mostly academic papers, standards, and vendor documentation.
I'm a bit unsure about excessive use of web resources (adacomputerscience.org, DigitalOcean, Modal blog, Stack Overflow)
textbook references would be stronger... if avaiable.
From my understanding, key papers are covered: Markidis et al.~[5], Fasi et al.~[21], Mukunoki et al.~[19], Ozaki et al.~[52,53], Higham [18].

\section{Strengths, Weaknesses, Structure}

\textbf{Strengths:}
\begin{itemize}
  \item Rigorous and reproducible (command-line args, seeds, hardware specs, software versions, ...). Everything is explicit.
  \item Theoretical error bound is derived first and experimentally validated.
  \item Honest reporting of negative results (overhead for small matrices, the \q{$e_{\text{avg}}$ trend when using 1 FP16 matrix} anomaly is openly investigated in Section 8.3).
  \item Adaptive splitting is a genuinely useful and novel engineering optimisation.
  \item Explicit acknowledgement of AI tool usage.
\end{itemize}

\textbf{Weaknesses}:
\begin{itemize}
  \item The main weakness: no empirical comparison with closest competitor (Ozaki scheme).
  \item Arguably narrow novelty, but absolutely fine in the context of a thesis.
  \item As argued above: evaluation uses only synthetic uniformly-random inputs
  \item Sustainability claims (SDG 9, 12) are unsubstantiated --- energy is never measured.
  \item Threshold (NZT=0.25) is fixed without justification or sensitivity analysis.
\end{itemize}

Section 4.5 (error bound derivation) is dense and hard for the reader.
It's the most difficult to understand.
I struggle to give suggestions on how to write it better, though.

Structurally, is standard and easy to follow.
Chapters are mostly well sized in relation.
The author has a tendency to be verbally extensive.
The background (Ch.~2) is simply too long.
Discussion (Ch.~8) is mostly recap of Ch.~7 --- in my opinion it could be sharper.

\section{Outlook and Concluding Assessment}
The work is nothing but solid.
The three-precision composition is novel as an explicit framing, however, mechanically close to prior work.
Adaptive splitting is a useful, if narrow, contribution.

Ethics are handled adequately.
Data is synthetic, sources are credited.
Sustainability and Development Goals Nr. 12 (Responsible Consumption) is claimed but \emph{not} substantiated
--- no energy or CO2 estimations.
This is a missed opportunity, especially given Future Work suggests it, because main findings might help improve computational efficiency.

A competent, well-executed engineering thesis with rigorous error analysis and honest experimental reporting.
The contribution is an extension of precision refinement to three precisions plus a simple runtime filter.
It is clearly motivated, well-bounded, and the empirical investigation is thorough within chosen scope.
Main weaknesses are the limited evaluation breadth (one GPU, synthetic inputs, no competitor benchmarks).

\section{Seminar Questions}
\begin{enumerate}[leftmargin=*]
  \item Your recommendation is \q{use 3 FP16 matrices}. Does/How does this recommendation change for real-world inputs?
    E.g.~covariance matrices or density matrices may have values that are heavily skewed rather than uniform on $[a,b]$?
  \item The NZT threshold is fixed at 0.25. How sensitive are $t_{\text{total}}$ and $e_{\text{avg}}$ to this choice? Is there a way to pick NZT from properties of the input matrix?
  \item Memory consumption prevents you from running $n=32768$.
    Could the intermediate $A_iA_j$ products be streamed/fused into the combination step to cut peak memory, given that the combination kernel sums them anyway?
  \item You made certain assumptions for your theoretical error bound. What assumption could you drop to tighten it and how would that change the bound?
  \item From your perspective, what is missing in your thesis? How could further improvement on precision refinement look like and what alternative approaches look promising?
  \item On what basis did you mention the increase of energy efficiency of your work?
    In what ballpark could the savings be?
\end{enumerate}

\end{document}
